\section{Requirement Check}

The requirement was checked at Mach \(0.9\), altitude \(10~\text{m}\), using the nominal sustainer point. The frequency-response requirement was evaluated using closed-loop bandwidth, because it directly measures the usable response frequency of each autopilot.

\begin{table}[H]
\centering
\caption{Requirement verification.}
\begin{tabular}{lccc}
\toprule
Requirement & Result & Limit & Status \\
\midrule
Lateral acceleration capability & \(13.52g\) & \(\geq 5g\) & Satisfied \\
Attitude autopilot bandwidth & \(1.652~\text{Hz}\) & \(\geq 0.5~\text{Hz}\) & Satisfied \\
Acceleration autopilot bandwidth & \(0.537~\text{Hz}\) & \(\geq 0.5~\text{Hz}\) & Satisfied \\
\bottomrule
\end{tabular}
\end{table}

The lateral acceleration capability was estimated from the acceleration transfer function and the \(10^\circ\) fin-deflection limit. The result is above \(5g\), so the missile has enough control authority at the design condition.

The acceleration autopilot is the limiting case. The original root-locus point gave an external gain that was too aggressive for the nominal model, so \(K_a\) was reduced to \(0.025\). This keeps the loop stable while preserving a bandwidth slightly above \(0.5~\text{Hz}\). The response is therefore acceptable for the stated requirements, although it is more oscillatory than the attitude loop.
